\chapter{Arzelà--Ascoli 定理}
\label{chap:arzela-ascoli}

$C(K)$ 的闭单位球通常不紧，是因为范数只限制函数值的大小，并不限制函数
在相近点之间怎样变化。一个函数族即使全部落在 $[-1,1]$ 内，也可以出现
越来越尖的峰或越来越快的振荡。Arzelà--Ascoli 定理表明，在紧致定义域上，
只需再加入对这种局部变化的一致控制，便能恢复有限维空间中有界序列所具有
的子列紧致性。

本章设 $(K,d)$ 为非空紧致度量空间，$m\ge1$ 为整数，
$\mathcal F\subset C(K;\R^m)$，值域采用欧氏范数，函数范数为
$\norm f_\infty=\max_{x\in K}\norm{f(x)}$。空函数族的等度连续、
有界及相对紧致结论均成立，共同连续模约定为零；以下证明只需考虑非空族。
若定义域为空，连续映射空间只有空映射，紧致性结论也直接成立。

\section{等度连续与一致有界}

\begin{definition}
\label{def:equicontinuity}
函数族 $\mathcal F$ 称为在 $x_0\in K$ 等度连续，如果对每个
$\varepsilon>0$，存在 $\delta>0$，使
\[
 d(x,x_0)<\delta\quad\Longrightarrow\quad
 \norm{f(x)-f(x_0)}<\varepsilon
\]
对所有 $f\in\mathcal F$ 同时成立。若它在每一点等度连续，则称
$\mathcal F$ 等度连续。
\end{definition}

这里的 $\delta$ 可以依赖于 $x_0$ 和 $\varepsilon$，但不能依赖于族中的
函数。有限个连续函数组成的族总是等度连续；无限族则未必。例如
$f_n(x)=x^n$ 在 $[0,1]$ 上不是等度连续的，因为在端点 $1$ 附近的变化越来越
集中。

紧致性把逐点的等度连续提升为全局的一致控制。

\begin{proposition}
\label{prop:equicontinuity-uniform-on-compact}
若 $\mathcal F$ 在紧致度量空间 $K$ 上等度连续，则对每个
$\varepsilon>0$，存在 $\delta>0$，使
\[
 d(x,y)<\delta\quad\Longrightarrow\quad
 \norm{f(x)-f(y)}<\varepsilon
\]
对所有 $f\in\mathcal F$ 同时成立。
\end{proposition}

\begin{proof}
对每个 $a\in K$ 选 $r_a>0$，使 $d(x,a)<2r_a$ 时
$\norm{f(x)-f(a)}<\varepsilon/2$ 对全部 $f$ 成立。有限个球
$B(a_j,r_{a_j})$ 覆盖 $K$。取 $\delta$ 为这些半径的最小值。若
$d(x,y)<\delta$，选 $j$ 使 $x\in B(a_j,r_{a_j})$；则 $x,y$ 都在
$B(a_j,2r_{a_j})$ 内，由三角不等式得到所需估计。
\end{proof}

函数族称为逐点有界，如果每个 $x\in K$ 都有
$\sup_{f\in\mathcal F}\norm{f(x)}<\infty$；称为一致有界，如果存在一个
$M$ 使 $\norm f_\infty\le M$ 对所有 $f$ 成立。一般而言逐点有界远弱于
一致有界，但与等度连续结合后，在紧集上两者相同。

\begin{proposition}
\label{prop:equicont-pointwise-to-uniform-bounded}
紧致 $K$ 上的等度连续、逐点有界函数族必一致有界。
\end{proposition}

\begin{proof}
在命题 \ref{prop:equicontinuity-uniform-on-compact} 中取
$\varepsilon=1$，得到 $\delta$。由紧致性选有限 $\delta$-网
$x_1,\ldots,x_N$。逐点有界使
$M_0=\max_j\sup_{f\in\mathcal F}\norm{f(x_j)}<\infty$。任给 $x$，选
$x_j$ 使 $d(x,x_j)<\delta$，便有
$\norm{f(x)}\le\norm{f(x_j)}+1\le M_0+1$。
\end{proof}

把两条命题合在一起，可以用一个函数记录整族的局部变化。若
$\mathcal F$ 等度连续且逐点有界，对 $t\ge0$ 定义
\[
 \omega_{\mathcal F}(t)
 =\sup\{\norm{f(x)-f(y)}:f\in\mathcal F,\ d(x,y)\le t\}.
\]
一致有界性保证这个上确界有限；它随 $t$ 单调不减，而紧致性上的共同
$\delta$ 说明 $\omega_{\mathcal F}(t)\to0$ 当 $t\downarrow0$。单个函数的
连续模因此被一个共同连续模取代。Lipschitz 族满足
$\omega_{\mathcal F}(t)\le Lt$，Hölder 族满足
$\omega_{\mathcal F}(t)\le Ct^\alpha$。

\section{稠密点上的对角抽取}

紧致度量空间是可分的。事实上，对每个 $n\ge1$ 取有限 $1/n$-网，把这些有限
集合求可数并，便得到至多可数稠密集 $D$。可写成
$D=\{x_1,x_2,\ldots\}$；若 $D$ 有限，则重复枚举其中的点。

设 $(f_n)$ 来自逐点有界族。数列 $(f_n(x_1))$ 在有限维空间中有收敛子列；
从这条子列中再抽取，使 $x_2$ 处的值收敛；继续下去。取第 $j$ 层子列的
第 $j$ 项组成对角子列 $(g_j)$，则对每个固定 $k$，$(g_j(x_k))$ 最终落在
第 $k$ 层子列中，因而收敛。

只在稠密集上收敛还不能保证一致收敛。等度连续性补上了稠密点之间的空隙。

\begin{lemma}[稠密点收敛与一致 Cauchy 控制]
\label{lem:dense-convergence-equicont-uniform-cauchy}
若 $(g_j)$ 等度连续，并且在稠密集 $D\subset K$ 的每一点都收敛，则
$(g_j)$ 在上确界范数下为 Cauchy 列。
\end{lemma}

\begin{proof}
给定 $\varepsilon>0$，由一致等度连续性取 $\delta>0$，使
$d(x,y)<\delta$ 时
$\norm{g_j(x)-g_j(y)}<\varepsilon/3$ 对每个 $j$ 成立。从 $D$ 中选有限
$\delta$-网 $x_1,\ldots,x_N$。每个 $(g_j(x_i))$ 都是 Cauchy 列，有限取
最大可得 $J$，使 $j,k\ge J$ 时这些网点上的差均小于 $\varepsilon/3$。
任给 $x$，选相邻网点 $x_i$，则
\[
 \norm{g_j(x)-g_k(x)}
 \le\norm{g_j(x)-g_j(x_i)}+
    \norm{g_j(x_i)-g_k(x_i)}+
    \norm{g_k(x_i)-g_k(x)}<\varepsilon.
\]
取上确界即得结论。
\end{proof}

\section{Arzelà--Ascoli 定理}

\begin{theorem}[Arzelà--Ascoli]
\label{thm:arzela-ascoli-sequential}
设 $K$ 为非空紧致度量空间，$m\ge1$。若序列 $(f_n)\subset C(K;\R^m)$ 等度连续且逐点
有界，则它有一致收敛子列。
\end{theorem}

\begin{proof}
取 $K$ 的至多可数稠密集。由逐点有界性和有限维 Bolzano--Weierstrass 定理，
按上一节作对角抽取，得到在稠密集每一点都收敛的子列 $(g_j)$。引理
\ref{lem:dense-convergence-equicont-uniform-cauchy} 说明它在 $C(K;\R^m)$
中为 Cauchy 列。每个坐标 $g_j^{(r)}$ 因而在标量空间 $C(K)$ 中为
Cauchy 列。由定理 \ref{thm:ck-banach}，它一致收敛于连续函数 $g^{(r)}$。
令 $g=(g^{(1)},\ldots,g^{(m)})$，则 $g$ 连续，且
\[
 \norm{g_j-g}_\infty
 \le\sqrt m\max_{1\le r\le m}\norm{g_j^{(r)}-g^{(r)}}_\infty
 \longrightarrow0.
\]
这也证明了这里使用的向量值连续函数空间的完备性。
\end{proof}

逐点有界可以换成一致有界，后者在应用中往往更容易验证。反过来，若一个
函数族的闭包在 $C(K;\R^m)$ 中紧致，则它必逐点有界且等度连续。

\begin{proof}[必要性的说明]
紧集在连续映射 $f\mapsto f(x)$ 下的像紧，故每一点的函数值有界。若闭包
不等度连续，则存在 $\varepsilon_0>0$、$f_n$ 和 $x_n,y_n$，使
$d(x_n,y_n)\to0$ 而
$\norm{f_n(x_n)-f_n(y_n)}\ge\varepsilon_0$。由紧致性抽取
$f_{n_k}\to f$ 一致；再由 $K$ 的紧致性抽取 $x_{n_k}\to x$，于是也有
$y_{n_k}\to x$。一致收敛和 $f$ 的连续性使上述差趋于零，矛盾。
\end{proof}

由此得到集合形式的完整刻画。

\begin{theorem}[相对紧致性判据]
\label{thm:arzela-ascoli-relative-compactness}
对 $\mathcal F\subset C(K;\R^m)$，下列条件等价：
\begin{enumerate}
 \item $\overline{\mathcal F}$ 在上确界范数下紧致；
 \item $\mathcal F$ 等度连续且逐点有界；
 \item $\mathcal F$ 等度连续且一致有界。
\end{enumerate}
\end{theorem}

\begin{proof}
前面的必要性给出 (1)$\Rightarrow$(2)，命题
\ref{prop:equicont-pointwise-to-uniform-bounded} 给出
(2)$\Rightarrow$(3)。设 (3) 成立，并取任意序列
$(g_n)_{n\ge1}\subset\overline{\mathcal F}$。对每个 $n\ge1$ 选
$f_n\in\mathcal F$，使 $\norm{g_n-f_n}_\infty<1/n$。由定理
\ref{thm:arzela-ascoli-sequential}，某个子列 $f_{n_k}$ 一致收敛于
$f\in C(K;\R^m)$；于是 $g_{n_k}$ 也一致收敛于 $f$。这证明闭包序列紧致。
度量空间中序列紧致等价于紧致，故 (3)$\Rightarrow$(1)。
\end{proof}

条件中的两个部分彼此独立。以下 $n\ge1$；常函数族 $f_n\equiv n$ 等度连续却不逐点有界；
$f_n(x)=\sin(nx)$ 在 $[0,2\pi]$ 上一致有界却不等度连续。常函数族的任何
子列在固定点都无有限极限。正弦函数列的任意子列 $\sin(n_kx)$ 也不能一致
收敛：若一致收敛于连续函数 $g$，则 $g(0)=0$。令
$x_k=\pi/(2n_k)\to0$，有
\[
 1=|\sin(n_kx_k)|
 \le\norm{\sin(n_k\cdot)-g}_\infty+|g(x_k)|\longrightarrow0,
\]
矛盾。这一论证处理了每个子列，而不仅是原函数族。

\section{两个常用判据}

\begin{corollary}[导数有界的函数族]
\label{cor:bounded-derivative-ascoli}
设 $a<b$、$\mathcal F\subset C^1([a,b];\R^m)$。若对某个 $x_0\in[a,b]$，
$\{f(x_0):f\in\mathcal F\}$ 有界，且
$\sup_{f\in\mathcal F}\norm{f'}_\infty<\infty$，则 $\mathcal F$ 在
$C([a,b];\R^m)$ 中相对紧致。
\end{corollary}

\begin{proof}
中值不等式给出共同的 Lipschitz 常数，故等度连续；再由
$\norm{f(x)}\le\norm{f(x_0)}+M|x-x_0|$ 得一致有界。应用相对紧致性判据。
\end{proof}

积分算子提供另一类重要例子。设 $a<b$、$c<d$，
$k\in C([a,b]\times[c,d];\R)$，对 $f\in C([c,d])$ 定义
\[
 (Tf)(x)=\int_c^d k(x,y)f(y)\,dy.
\]
由积分的线性性，$T$ 是线性的，且
$\norm{Tf}_\infty\le(d-c)\norm k_\infty\norm f_\infty$。
核的一致连续性保证 $Tf$ 连续，故 $T:C([c,d])\to C([a,b])$ 是有界线性算子。
对 $C([c,d])$ 的单位球，像函数一致有界；而
\[
 |Tf(x)-Tf(x')|
 \le(d-c)\sup_y|k(x,y)-k(x',y)|.
\]
$k$ 的一致连续性使右端与 $f$ 无关地趋于零。故 $T$ 把有界集送入相对紧
的函数族。把有界集映成相对紧集的有界线性算子称为\emph{紧算子}。
这里核的共同连续模使单位球的像满足 Arzelà--Ascoli 的条件，尽管定义域
与值域本身都是无限维空间。

\section*{习题}
\addcontentsline{toc}{section}{习题}

\subsection*{A组　等度连续}
\begin{enumerate}[label=\textbf{\thechapter.\arabic*.}]
 \exerciseitem{finite-family-equicontinuous}
 证明紧致空间上的有限连续函数族等度连续。
 \exerciseitem{power-family-not-equicontinuous}
 直接从定义证明 $\{x^n:n\ge1\}$ 在 $[0,1]$ 上不等度连续。
 \exerciseitem{equicontinuity-uniform-proof-details}
 检查命题 \ref{prop:equicontinuity-uniform-on-compact} 中球半径的选择，并给出
 不紧定义域上结论失败的例子。
 \exerciseitem{pointwise-bounded-to-uniform}
 不引用正文证明，重新证明等度连续且逐点有界的函数族在紧集上一致有界。
 \exerciseitem{common-modulus-equicontinuity}
 设函数族等度连续且逐点有界。证明存在递增函数 $\omega$，当
 $t\downarrow0$ 时有 $\omega(t)\downarrow0$，并且
 \[
  \norm{f(x)-f(y)}\le\omega(d(x,y)).
 \]
\end{enumerate}

\subsection*{B组　抽取与紧致性}
\begin{enumerate}[label=\textbf{\thechapter.\arabic*.},resume]
 \exerciseitem{compact-metric-separable}
 详细证明每个紧致度量空间都有至多可数稠密子集。
 \exerciseitem{diagonal-subsequence-bookkeeping}
 写出对角子列的严格指标构造，并证明它在每个稠密点收敛。
 \exerciseitem{dense-convergence-not-uniform-without-equicont}
 构造在稠密集上收敛但没有一致收敛子列的连续函数序列。
 \exerciseitem{ascoli-necessity-direct}
 补全相对紧族必等度连续的反证法，明确每次抽取所用的紧致性。
 \exerciseitem{ascoli-closed-family-compact}
 证明等度连续、逐点有界且在上确界范数下闭的函数族是紧集。
\end{enumerate}

\subsection*{C组　应用}
\begin{enumerate}[label=\textbf{\thechapter.\arabic*.},resume]
 \exerciseitem{bounded-derivative-family-compactness}
 证明推论 \ref{cor:bounded-derivative-ascoli}，并说明只控制导数而不控制一个
 点的函数值为何不够。
 \exerciseitem{holder-family-relative-compact}
 设 $M,C\ge0$、$\alpha>0$ 为固定常数。证明满足
 $|f(x)|\le M$、$|f(x)-f(y)|\le C|x-y|^\alpha$ 的函数族在
 $C([0,1])$ 中相对紧致。
 \exerciseitem{integral-operator-compact}
 详细证明连续核积分算子把任意有界集送入相对紧集。
 \exerciseitem{primitive-family-ascoli}
 若 $f_n\in C([0,1])$ 且 $|f_n|\le1$，令
 $F_n(x)=\int_0^xf_n(t)\,dt$。证明 $(F_n)$ 有一致收敛子列，并说明不能
 由此断言 $(f_n)$ 有一致收敛子列。
 \optionalexerciseitem{ascoli-metric-valued}
 把定理推广到取值于完备度量空间 $Y$ 的连续映射族，其中“逐点有界”应改为
 每个点值集合在 $Y$ 中相对紧。
\end{enumerate}
